


\begin{frame}[fragile]{Let, Var, and Const}

\begin{itemize}
	\item JavaScript has two ways to make a variable: let and var. 
	\item The difference between the two is that the variables defined with `let' cannot be redefined, whereas `var' variables are OK with the redefine. 
    \item The variables defined with `const' cannot be modified. 
\end{itemize}

\end{frame}

\begin{frame}[fragile]{}

\begin{itemize}
	\item The code shows an example. Line 6 is an error as we try to redefine the variable a that is defined with `let.' 
	\item It is OK to update the value defined with let as in line 7, or redefine the variable defined with var as in line 8. 
	\item It is not OK to change constant value as in line 9. 
\end{itemize}

\begin{Code}{}%
\begin{lstlisting}[firstnumber=1, label=glabels, xleftmargin=0pt, basicstyle=\tiny]% 

if (true) { // to make local variables
  let a = 1; 
  var b = 2;
  const c = 3;
  if (true) { // to make local variables
    var a = 4; // Error
    a = 3; // OK 
    let b = 3;  // OK
    c = 4; //Error 
  }
  console.log(a); 
  console.log(b);
}
\end{lstlisting}%   
\end{Code}

\end{frame}

\begin{frame}[fragile]{JavaScript String}

\begin{itemize}
	\item Java has a char type (primitive value) and a String (reference type object), and we use '.' to represent a char, and ''...'' for a Java String. 
	\item JavaScript has only strings, and we use one notation \verb|'...'|.
	\item We can concatenate strings and variables using the '+' operator as in line 2. 
	\item We also can interpolate strings (use variables in a string) with \verb|`...`| and \$\{...\} as in line 3.
\end{itemize}

\begin{Code}{}%
\begin{lstlisting}[firstnumber=1, label=glabels, xleftmargin=0pt, basicstyle=\scriptsize]% 

var cruel = 'cruel';
var line = 'hello ' + cruel + ' world';
var line2 = `hello ${cruel} world`;
\end{lstlisting}%   
\end{Code}%

\end{frame}

\begin{frame}[fragile]{}

\begin{itemize}
	\item When the interpolate string is used as an argument to a function, the function can analyze the input to extract values. 
	\item We call this feature `tagged list.'
	\item In this example, the analyzer() function parse the string input to extract var1 and var2. 
	\item Notice that this code is output from Quokka extension, so it displays parsed input results in lines 87 - 89. 
\end{itemize}

\begin{center}
\includegraphics[width=0.7\textwidth]{topics/SE_programming/javascript/pic/history/string2.png}
\end{center} 

\end{frame}